\section{Variáveis Aleatórias}

Variáveis aleatórias (v.a) são representações de características numéricas de
experimentos aleatórios. Essa visão extende-se desde dos casos mais simples, a exemplo de
quando o próprio resultado é um atributo numérico, quanto a características mais complexas
como a soma dos lados de dados, quantas vezes lançamos uma moeda até
obtermos duas caras ou o número de peças defeituosas produzidas após um dia de trabalho.
As v.a's podem ser formalizadas como funções que  associam cada resultado do
espaço amostral a um número real.

\begin{def:va}
Se $\Omega$ é um espaço amostral, então uma variável aleatória $X$ é uma aplicação
de $\Omega$ em $\mathbb{R}$.
suporte.
\end{def:va}

Com isso, podemos expressar eventos de um espaço amostral como aqueles que levam a variável
a ter uma realização específica. Por exemplo, se o experimento consiste em lançar dois
dados, e $X$ é a soma dos dados, então o evento dos resultados nos quais a soma do lados
sorteados é igual a 7 pode ser expressa como simplesmente $X = 7$. A proposição a seguir
justifica esse racioncínio.

\begin{prop:prob va}
Seja $X: \Omega \to \mathbb{R}$ uma v.a. discreta. Se $E \subset \Omega$ é tal que
\[
	E = \{\omega \in \Omega \mid X(\omega) = x\},
\] então
\[
	P(X = x) = P(E)
\]
\begin{proof}
	Temos que $X(\omega) = x$ se, e somente se, $\omega \in E$. Logo, os eventos $X = x$
	e $E$ são equivalentes e possuem, pois, a mesma probabilidade.
\end{proof}
\end{prop:prob va}

As variáveis tem associadas a si uma \textbf{função de probabilidade}. Em síntese,
uma função de probabilidade nos dá uma ideia do quão frequente é um valor ou intervalos
de valores de uma variável aleatória, além de poder informar sobre valores centrais e
variabilidade da variável. A interpretação bayesiana para essas funções seria como a nossa
crença modela a confiança \textit{a priori} em cada resultado -- ou intervalo deles --
possível do experimento. Uma família de funções de probabilidades, cujos membros
diferenciam-se a menos do valor de um ou mais parâmetros, são chamadas de
\textbf{distribuições de probabilidade}.
